\documentclass[11pt]{article}

\usepackage[margin=1in]{geometry}
\usepackage{amsmath, amssymb}
\usepackage{longtable}
\usepackage{array}

\newcommand{\toprule}{\hline}
\newcommand{\midrule}{\hline}
\newcommand{\bottomrule}{\hline}
\newcommand{\href}[2]{\texttt{\detokenize{#2}}}
\newcommand{\url}[1]{\texttt{\detokenize{#1}}}

\title{Adaptive Fractionation at the MR-Linac}
\author{
  Y. P\'erez Haas\thanks{Corresponding author: \href{mailto:yoelhaas@gmail.com}{yoelhaas@gmail.com}} \and
  R. Ludwig \and
  R. Dal Bello \and
  S. Tanadini-Lang \and
  J. Unkelbach
}
\date{
  Department of Radiation Oncology, University Hospital of Zurich, Zurich, Switzerland\\
  \vspace{0.5em}
  Received June 22, 2022; revised November 17, 2022; accepted January 3, 2023; published January 19, 2023
}

\begin{document}
\maketitle

\begin{center}
\textbf{Internal LaTeX transcription for repository reference}\\
Physics in Medicine \& Biology 68 (2023) 035003\\
DOI: \href{https://doi.org/10.1088/1361-6560/acafd4}{10.1088/1361-6560/acafd4}
\end{center}

\begin{abstract}
\textbf{Objective.}
Fractionated radiotherapy typically delivers the same dose in each fraction. Adaptive fractionation (AF) is an approach to exploit inter-fraction motion by increasing the dose on days when the distance of tumor and dose-limiting organs at risk (OAR) is large and decreasing the dose on unfavorable days. We develop an AF algorithm and evaluate the concept for patients with abdominal tumors previously treated at the MR-linac in 5 fractions.

\textbf{Approach.}
Given daily adapted treatment plans, inter-fractional changes are quantified by sparing factors $\delta_t$ defined as the OAR-to-tumor dose ratio. The key problem of AF is to decide on the dose to deliver in fraction $t$, given $\delta_t$ and the dose delivered in previous fractions, but not knowing future $\delta_t$ values. Optimal doses that maximize the expected biologically effective dose in the tumor (BED10) while staying below a maximum OAR BED3 constraint are computed using dynamic programming, assuming a normal distribution over $\delta$ with mean and variance estimated from previously observed patient-specific $\delta_t$ values. The algorithm is evaluated for 16 MR-linac patients in whom tumor dose was compromised due to proximity of bowel, stomach, or duodenum.

\textbf{Main results.}
In 14 out of the 16 patients, AF increased the tumor BED10 compared to the reference treatment that delivers the same OAR dose in each fraction. However, in 11 of these 14 patients, the increase in BED10 was below 1 Gy. Two patients with large sparing factor variation had a benefit of more than 10 Gy BED10 increase. For one patient, AF led to a 5 Gy BED10 decrease due to an unfavorable order of sparing factors.

\textbf{Significance.}
On average, AF provided only a small increase in tumor BED. However, AF may yield substantial benefits for individual patients with large variations in the geometry.
\end{abstract}

\noindent\textbf{Keywords:} adaptive fractionation, adaptive radiotherapy, MR-linac, MR-guided radiotherapy, inter-fraction motion

\section{Introduction}

Most radiation treatments are fractionated because normal tissue can tolerate higher doses if the radiation dose is split into several fractions (Lajtha et al. 1960; Fowler 2006). Motion of tumors and organs at risk (OAR) in between fractions is generally assumed to degrade the quality of treatments. Traditionally, safety margins are used to account for motion, which worsens the trade-off between tumor coverage and normal tissue sparing that would be possible without motion (van Herk et al. 2000).

Nowadays, image guidance technology allows for measuring organ motion during the treatment, and the treatment plan can be adapted to the changing anatomy (Mayinger et al. 2021). Thereby, safety margins can be reduced and, in the ideal case, such adaptive radiotherapy concepts restore the treatment quality that is possible for a static patient without motion (Guckenberger et al. 2011; Chen et al. 2013; Brock 2019). This led to a widespread implementation of stereotactic body radiation (SBRT) (Wulf et al. 2006; Lo et al. 2010; Andratschke et al. 2018).

However, in principle inter-fraction tumor motion can even be exploited, that is, a better treatment quality may be achieved in the presence of motion compared to a static patient geometry. Adaptive fractionation (AF) (Chen et al. 2008; Lu et al. 2008; Ramakrishnan et al. 2012) is one approach to exploit inter-fraction motion. In this approach, a treatment plan is not only adapted to the daily patient geometry, but also the dose delivered to the tumor in each fraction is modified: the dose is increased on favorable treatment days, i.e. when the distance between tumor and dose-limiting OAR is relatively large; and the dose is reduced for unfavorable geometries, i.e. when the tumor and OAR are closer. Thereby, the ratio between total dose delivered to the OAR versus total dose delivered to the tumor may be improved compared to uniform treatments that deliver the same dose in each fraction.

Although the idea of AF has been presented previously, clinical translation has been facing substantial hurdles (Chen et al. 2008; Lu et al. 2008; Ramakrishnan et al. 2012). AF relies on imaging to measure the changes of the patient geometry from day to day (Wu et al. 2002). This became possible to some degree with cone-beam CT. However, the limited soft tissue contrast limits the potential clinical applications. Magnetic-resonance (MR) guided radiotherapy (Acharya et al. 2016; Henke et al. 2018; Palacios et al. 2018; Kl\"uter 2019) allows the visualization of anatomical variation with higher soft tissue contrast. This extends the range of potential applications to abdominal lesions in proximity to bowel and stomach, which can exhibit very large inter-fraction motion and may thus benefit from AF.

In this work, an approach for AF at an MR-linac is presented and evaluated. The main contributions of this paper are:

\begin{enumerate}
\item To compute the optimal dose to deliver in each fraction, knowing today's geometry and the dose delivered in previous fractions, but not knowing the patient geometry in future fractions, a dynamic programming algorithm has been developed, extending the work of Ramakrishnan et al. (2012).
\item The algorithm is tested on patients previously treated at the MR-linac with 5-fraction SBRT for abdominal lesions near bowel, stomach, or duodenum. Thus, the potential benefit of AF is estimated for real patient data, extending previous work that only conceptually introduced AF based on synthetic data (Chen et al. 2008; Lu et al. 2008; Ramakrishnan et al. 2012).
\end{enumerate}

\section{Methods and Materials}

\subsection{Patients and treatment plans}

We consider patients with abdominal tumors in proximity to either the bowel, stomach, or duodenum who received 5-fraction SBRT treatments at the MR-linac system (MRIdian, ViewRay). In all cases, tumor coverage was compromised due to the dose received by the dose-limiting OAR. All patients were planned and treated according to institutional practice. In addition to the simulation MR and CT scans, daily MR scans were performed for online adaptive radiotherapy. Tumors and OARs in a 2 cm ring around the tumor were recontoured according to institutional guidelines and daily adaptive treatment plans were created. Thus, the dose distributions were reoptimized in each fraction to adapt to inter-fractional changes, without altering the prescription dose.\footnote{GTV and PTV prescription doses and OAR constraints varied between patients. However, for all patients coverage was compromised due to the OAR constraint.}

For the purpose of this project, dose-volume histograms (DVHs) of 16 patients were exported from the treatment planning system. DVHs were exported for GTV, PTV, and the relevant OARs for 6 treatment plans per patient, corresponding to the 5 delivered plans and the initial plan based on the planning MR.

\subsection{Sparing factors}

For the purpose of AF, treatment plans and the geometric variations are described in terms of sparing factors $\delta$:

\begin{equation}
\delta_t = \frac{d_t^N}{d_t},
\label{eq:sparing-factor}
\end{equation}

where $d_t^N$ denotes the dose received by the dose-limiting OAR in fraction $t$ and $d_t$ is the dose received by the tumor. For the definition of $d_t^N$ and $d_t$, we follow the clinical practice of dose prescription and constraint specification: the dose to the OAR $d_t^N$ is defined as the dose exceeded in 1 cc of the OAR ($D_{1\mathrm{cc}}$), which represents a commonly used dose parameter for the bowel, stomach, or duodenum in SBRT treatments. The tumor dose $d_t$ is defined as the dose exceeded in 95\% of the PTV volume ($D_{95}$), which is the dose parameter commonly used for dose prescription and reporting (ICRU 1993, 1999). Each patient is thus described via a sequence of 6 sparing factors corresponding to the planning MR and the 5 fractions. We further assume that inter-fraction motion is random and is described by a Gaussian distribution over $\delta$ with a patient-specific mean $\mu$ and standard deviation $\sigma$:

\begin{equation}
\delta_t \sim \mathcal{N}(\mu, \sigma^2).
\label{eq:delta-distribution}
\end{equation}

\subsection{Fractionation}

To model the fractionation effect, the biologically effective dose (BED) model is used (Jones et al. 2001; McMahon 2018). It is assumed that the classic BED model can be extended to varying doses per fraction such that the cumulative BED at the end of treatment is given by the sum of the BED values delivered in individual fractions. Thus, the cumulative BED delivered to the tumor is given by

\begin{equation}
B_T = \sum_{\tau=1}^{t}\left(d_{\tau} + \frac{d_{\tau}^2}{(\alpha/\beta)_T}\right),
\label{eq:bed-tumor}
\end{equation}

where $d_{\tau}$ denotes the dose delivered to the tumor in fraction $\tau$. Consequently, the cumulative BED received by the OAR is

\begin{equation}
B_N = \sum_{\tau=1}^{t}\left(\delta_{\tau} d_{\tau} + \frac{\delta_{\tau}^2 d_{\tau}^2}{(\alpha/\beta)_N}\right),
\label{eq:bed-oar}
\end{equation}

where $\delta_{\tau}$ is the sparing factor in fraction $\tau$. In this work, the $\alpha/\beta$ ratios for the OARs and the tumors are the same for all patients and set to $(\alpha/\beta)_N = 3$ and $(\alpha/\beta)_T = 10$ (van Leeuwen et al. 2018). Correspondingly, the biological effective doses in the OAR and the tumor will be denoted as BED3 and BED10 respectively.

Note that the calculation of cumulative BED3 in equation~\eqref{eq:bed-oar} assumes that the same 1 cc of the dose-limiting OAR receives the highest dose, which may not be the case in reality. In this case, equation~\eqref{eq:bed-oar} can be considered a worst-case measure for OAR dose, which overestimates the cumulative BED3 received by any part of the OAR. However, due to the impracticality of deformable dose accumulation in the abdomen, the same approximation is done in current clinical practice.

\subsection{Adaptive fractionation}

The goal of AF is to optimally decide on the doses $d_t$ delivered to the tumor in each fraction so as to maximize the expected cumulative BED10 delivered to the tumor, subject to a constraint on the cumulative OAR BED. For 5-fraction SBRT treatments in the abdomen, limiting the dose to bowel, stomach, and duodenum is prioritized, i.e. target coverage is compromised if necessary to fulfill the OAR constraint. This approach is clinical practice at our institution (Mayinger et al. 2021) as well as other clinics (Tyagi et al. 2021). For this work, we assume that the BED delivered to the dose-limiting OAR may not exceed $B_N^{\max} = 90$ Gy BED3, which corresponds to 30 Gy physical dose delivered in 5 uniform fractions (Pavic et al. 2022). Thus, for patients in whom the tumor dose is compromised due to the OAR constraint, as considered in this paper, the goal of AF is to increase tumor BED.\footnote{For applications of AF to other treatment sites where target coverage is prioritized over OAR sparing, the goal would instead be to minimize expected cumulative OAR BED subject to a constraint on delivering the prescribed cumulative tumor BED.}

In each fraction $t$, the sparing factor $\delta_t$ can be determined. In addition, we know the BED delivered to the tumor and the OAR in previous fractions. Based on this information, the decision must be made which dose should be delivered in today's fraction. The difficulty comes from not knowing whether the remaining future fractions will have favorable or unfavorable patient geometries. The sparing factors are random variables with an estimated probability distribution, but the exact future values are unknown.

From a practical perspective, this approach to AF would be implemented by up-scaling or down-scaling the reoptimized treatment plan for that fraction. That is, we assume that the adaptive radiotherapy process consisting of MR imaging, recontouring, and plan reoptimization is not altered. The only additional step would be a final up-scaling or down-scaling of the fluence without changing the shape of the dose distribution. This corresponds to a renormalization of the plan, which is in current practice done within a narrow range that can be extended to implement AF.

\subsection{MDP model}

To determine the optimal doses $d_t$, we apply the framework of Markov decision processes (MDP) and formulate AF as a stochastic optimal control problem. Here, we first describe the MDP model for a known probability distribution $P(\delta_t)$ and afterwards discuss how to estimate and update $P(\delta_t)$. Optimal control problems are described by states, actions, state transitions, and reward functions. In this application, these are given by:

\paragraph{State.}
In each fraction, the state of a patient's treatment is described by a two-dimensional vector $s = (\delta, B_N)$ that specifies today's sparing factor $\delta$ and the cumulative BED of the OAR that has been delivered so far in previous fractions. Thus, the state of a treatment in fraction $t$ for a patient with sparing factors $\{\delta_{\tau}\}_{\tau=1}^{t}$ treated with doses $\{d_{\tau}\}_{\tau=1}^{t-1}$ is

\begin{equation}
s_t = \left(\delta_t, \sum_{\tau=1}^{t-1}\left(\delta_{\tau} d_{\tau} + \frac{\delta_{\tau}^2 d_{\tau}^2}{(\alpha/\beta)_N}\right)\right).
\label{eq:state-two-dim}
\end{equation}

\paragraph{Action and policy.}
The actions correspond to the physical doses $d_t$ that are delivered to the tumor in a fraction. Thus, a policy specifies for each fraction $t$ and possible state of the treatment, the dose that should be delivered in this state. Tumor doses are constrained by a maximum dose per fraction $d_{\max}$ and a minimum dose per fraction $d_{\min}$. In the results below, we assume $d_{\min} = 0$ Gy and $d_{\max} = 23$ Gy. A single dose of 23 Gy allows delivery of the BED10 equivalent of a typical prescription dose of 40 Gy in 5 fractions in a single fraction. This range was chosen to explore the full potential of AF. The impact of tighter constraints on the dose per fraction, which may be considered in a clinical implementation, is discussed in section~4.1.

\paragraph{State transition.}
If, in fraction $t$, the treatment is in state $s_t = (\delta_t, B)$ and a dose $d_t$ is delivered to the tumor, the state transitions to

\begin{equation}
s_{t+1} = \left(\delta_{t+1}, B + \delta_t d_t + \frac{\delta_t^2 d_t^2}{(\alpha/\beta)_N}\right),
\label{eq:state-transition}
\end{equation}

in fraction $t+1$. The BED component of the future state is calculated by adding the OAR BED delivered in fraction $t$ to the previously delivered BED $B$, which is assumed deterministic (i.e. we do not consider uncertainty in dose delivery). The sparing factor in fraction $t+1$ is random, making the state transition probabilistic. The probability distribution for the state transition is simply given by the probability distribution over the sparing factors, $P(\delta)$.

\paragraph{Reward.}
In each fraction $t$, the immediate reward $r_t$ is given by the BED delivered to the tumor in that fraction:

\begin{equation}
r_t = d_t + \frac{d_t^2}{(\alpha/\beta)_T}.
\label{eq:reward}
\end{equation}

To account for the cumulative BED constraint in the dose-limiting OAR, the BED must be below $B_N^{\max}$. To enforce this, a terminal reward of $-\infty$ is assigned to all terminal states in which the cumulative OAR BED delivered after the last fraction exceeds the constraint value.

A characteristic of the specified MDP model is that the cumulative BED delivered to the tumor is not part of the state $s$. It is only integrated in the reward $r_t$. The reason for this is that the optimal policy does not depend on the tumor state, i.e. the optimal dose to deliver does not depend on the previously accumulated tumor BED. The intuitive explanation is that, in each state and fraction, we aim to maximize future tumor BED, independent of the previously accumulated tumor BED. However, if the goal was to deliver a fixed prescribed tumor BED by the end of the treatment, the cumulative tumor BED would be part of the state (see also section~4.3).

\subsection{Dynamic programming algorithm}

A dynamic programming (DP) algorithm can be used to compute the optimal policy with the help of a value function (Sutton and Barto 2018). The value function $v_t(\delta, B)$ describes how desirable it is to be in state $s_t(\delta, B)$ in fraction $t$ and, therefore, it contains the information whether an action should be taken to reach that state. In this application, the value for each state represents the expected cumulative BED that can be delivered to the tumor in the remaining fractions, starting from that state and acting according to the optimal policy.

The Bellman equation relates the value function in fraction $t$ to the optimal policy and the value function in the subsequent fraction, which for this application reads

\begin{equation}
v_t(\delta, B) = \max_{d}\left[d + \frac{d^2}{(\alpha/\beta)_T} + \sum_{\delta'} P(\delta')\, v_{t+1}\left(\delta', B + \delta d + \frac{\delta^2 d^2}{(\alpha/\beta)_N}\right)\right],
\label{eq:bellman}
\end{equation}

and

\begin{equation}
\pi_t(\delta, B) = \arg\max_{d}\left[d + \frac{d^2}{(\alpha/\beta)_T} + \sum_{\delta'} P(\delta')\, v_{t+1}\left(\delta', B + \delta d + \frac{\delta^2 d^2}{(\alpha/\beta)_N}\right)\right].
\label{eq:policy}
\end{equation}

Value function and optimal policy can be calculated iteratively in one backward recursion starting from the last fraction. To that end, the value function $v_{F+1}$, corresponding to the terminal reward at the end of the treatment after all $F$ fractions are delivered, is initialized to

\[
v_{F+1}(\delta, B) =
\begin{cases}
-\infty & \text{if } B > B_N^{\max}, \\
0 & \text{if } B \leq B_N^{\max}.
\end{cases}
\]

The optimal policy is found by discretizing both actions and states. To reduce discretization artifacts in the optimal policy, two tricks are applied. First, linear interpolation of the value function in the BED component of the state is used. Second, one can exploit that, in the optimal policy, the last fraction will simply deliver the maximum residual BED3 to the OAR to end up at the cumulative BED3 constraint $B_N^{\max}$. This can be used to directly initialize value function and optimal policy in the last fraction $F$ using continuous values of $d_F$. Thereby, artifacts from not reaching exactly $B_N^{\max}$ in the last fraction due to discretization of the actions can be avoided. The runtime of the algorithm is on the order of seconds and therefore suited for online treatment adaptation.

\subsection{Probability updating}

The DP algorithm relies on a description of the environment to compute an optimal policy, in this case the probability distribution of the sparing factor $P(\delta)$, which we assume to be a Gaussian distribution truncated at 0, with patient-specific parameters for mean and standard deviation. At the start of a treatment, only two sparing factors are available for that patient, from the planning scan and the first fraction. In each fraction, an additional sparing factor is measured, which can be used to calculate updated estimates $\mu_t$ and $\sigma_t$ for mean and standard deviation, respectively.

In each fraction $t$, we use the maximum likelihood estimator for the mean of the sparing factor distribution given by

\begin{equation}
\mu_t = \frac{1}{t+1}\sum_{\tau=0}^{t}\delta_{\tau},
\label{eq:mu-update}
\end{equation}

where $\delta_0$ denotes the sparing factor from the planning MR. The estimator for the standard deviation, given the patient-specific sparing factors up to fraction $t$, follows a chi-squared distribution, and the maximum likelihood estimator is given by

\begin{equation}
\sigma_t^{\mathrm{pat}} = \sqrt{\frac{1}{t+1}\sum_{\tau=0}^{t}\left(\delta_{\tau} - \mu_t\right)^2}.
\label{eq:sigma-pat}
\end{equation}

However, the standard deviation may be severely under- or overestimated if calculated from only two samples at the very beginning of the treatment. Therefore, we assume a population-based prior for the standard deviation and compute the maximum a posterior estimator of $\sigma_t$ via Bayesian inference. As the sparing factors are assumed to follow a normal distribution with unknown variance, a gamma distribution is chosen as prior to estimate the standard deviation $\sigma$:

\begin{equation}
f(\sigma; k, \theta) = \frac{1}{\Gamma(k)\theta^k}\sigma^{k-1}\exp\left(-\frac{\sigma}{\theta}\right),
\label{eq:gamma-prior}
\end{equation}

with shape parameter $k$ and scale parameter $\theta$. The maximum a posterior estimator for the standard deviation in fraction $t$ is then given by

\begin{equation}
\sigma_t = \arg\max_{\sigma}\left[\frac{\sigma^{k-1}}{\sigma^{t-1}} \exp\left(-\frac{\sigma}{\theta}\right)\exp\left(-\frac{(\sigma_t^{\mathrm{pat}})^2}{2\sigma^2/t}\right)\right].
\label{eq:sigma-map}
\end{equation}

Using equations~\eqref{eq:mu-update} and~\eqref{eq:sigma-map}, the probability distribution $P(\delta; \mu_t, \sigma_t)$ is updated with every newly acquired sparing factor and used in the Bellman equations~\eqref{eq:bellman} and~\eqref{eq:policy} to recompute the optimal policy before each fraction.

\subsection{Quantification of the benefit}

The treatment plan produced by AF is compared to three other treatments:

\begin{enumerate}
\item A reference treatment in which 6 Gy physical dose (18 Gy BED3) is delivered to the OAR in each fraction. Hence, the reference treatment delivers exactly the upper limit of 90 Gy BED3 to the OARs.
\item An upper bound for the benefit of AF. To do so, we consider the hypothetical situation that all sparing factors $\delta_t$ are known before treatment. In that case, the optimal doses per fraction $d_t$ are calculated by solving the following optimization problem:
\begin{equation}
\begin{aligned}
\text{maximize}_{d}\quad & \sum_{t=1}^{F}\left(d_t + \frac{d_t^2}{(\alpha/\beta)_T}\right) \\
\text{subject to}\quad & \sum_{t=1}^{F}\left(\delta_t d_t + \frac{\delta_t^2 d_t^2}{(\alpha/\beta)_N}\right) \leq B_N^{\max}.
\end{aligned}
\label{eq:upper-bound}
\end{equation}
This treatment would optimally exploit the variation in $\delta$ and can thus be used to benchmark the benefit of AF. However, it represents an unachievable upper bound for any realistic approach to AF where future sparing factors are unknown.
\item The clinically delivered treatment. The clinical treatment aims to deliver a fixed prescription dose to the tumor in each fraction and may deliver less than 90 Gy BED3 to the OAR. Hence, the clinical treatment is included for qualitative comparison, while the reference treatment is used for quantitative evaluation of the benefit of AF.
\end{enumerate}

To evaluate AF for a larger patient cohort, we generate additional patients by randomly sampling sparing factors from distributions that resemble the 16 patients previously treated at the MR-linac.

\begin{figure}[ht]
\centering
\fbox{\parbox{0.8\linewidth}{Figure omitted in this text transcription. See the original PDF for the rendered graphic.}}
\caption{Scatter plot of all acquired sparing factors.}
\label{fig:scatter-sparing-factors}
\end{figure}

\begin{figure}[ht]
\centering
\fbox{\parbox{0.8\linewidth}{Figure omitted in this text transcription. See the original PDF for the rendered graphic.}}
\caption{Gamma distribution for the population-based prior for the standard deviation (solid line), estimated from the observed standard deviations in the 16 patients (dotted blue lines).}
\label{fig:gamma-prior}
\end{figure}

\section{Results}

\subsection{Observed sparing factors}

Figure~\ref{fig:scatter-sparing-factors} shows the sparing factor distribution for each of the 16 patients and the respective dose-limiting OAR. Substantial inter-patient variation is observed regarding the intra-patient variation of the sparing factors. Some patients show substantial variation with a standard deviation of approximately 0.1 (e.g. patients 3, 5, 8, 13, and 16), whereas other patients show little variation of the sparing factor between fractions (e.g. patients 1, 2, and 9). The numerical values of the sparing factors as well as the DVHs from which they are computed are provided as CSV files in the supplementary materials. The DVHs of patients 7, 8, and 13 are shown in the appendix.

The standard deviations of all patients were used to compute the hyperparameters of the gamma prior. Figure~\ref{fig:gamma-prior} shows the resulting prior distribution, illustrating the inter-patient variation in the anatomical variability.

\subsection{Illustration of AF for an example patient}

To illustrate AF based on the DP algorithm, patient 7 is discussed in detail. The sparing factors observed are: $\delta_0 = 0.88$ for the planning MR and $[0.99, 0.87, 0.98, 1.04, 1.00]$ for the 5 fractions. In the first fraction, two sparing factors are known, $\delta_0$ and $\delta_1$. These two sparing factors lead to the probability distribution $P(\delta; \mu_1, \sigma_1)$ shown in figure~3(a) in the original paper, with a mean of $\mu_1 = 0.94$ and standard deviation $\sigma_1 = 0.058$.

\begin{table}[ht]
\centering
\caption{Dose delivered to the tumor for different treatments for patient 7. The results for AF in this table are computed with probability distribution updating. Thus, the doses are not identical to the ones depicted in figure 3 of the original paper.}
\label{tab:patient7}
\begin{tabular}{lllll}
\toprule
Fraction & Sparing factor & Upper bound & Adaptive fractionation & Reference plan \\
\midrule
Planning MR & 0.88 & --- & --- & --- \\
First fraction & 0.99 & 4.4 & 4.1 & 6.0 \\
Second fraction & 0.87 & 12.6 & 8.8 & 6.9 \\
Third fraction & 0.98 & 4.7 & 4.5 & 6.1 \\
Fourth fraction & 1.04 & 3.3 & 4.4 & 5.8 \\
Fifth fraction & 1.00 & 4.1 & 8.2 & 6.0 \\
Total BED tumor & --- & 51.7 & 50.2 & 50.0 \\
\bottomrule
\end{tabular}
\end{table}

Figure 3 in the original paper shows the optimal policy obtained for this probability distribution. In the first fraction, the OAR BED is zero and the dose to deliver depends only on the sparing factor $\delta_1$. Since $\delta_1$ is relatively high compared to the sparing factors to be expected given $P(\delta; \mu_1, \sigma_1)$, a low dose of $d_1 = 4.1$ Gy is delivered to the tumor. In the second fraction, the dose to be delivered depends on the OAR BED delivered in the first fraction $B$ and the sparing factor $\delta_2$. The optimal policy is monotone in $B$ and $\delta_2$, i.e. lower sparing factors and lower previously accumulated BED lead to higher doses delivered in the current fraction. Given the OAR BED3 of 9.6 delivered in the first fraction and the sparing factor $\delta_2 = 0.87$, the optimal dose to deliver in fraction 2 is $d_2 = 9.8$ Gy, a rather high dose as $\delta_2$ is significantly lower than the mean of the sparing factor distribution and a low dose was delivered in the first fraction.

The structure of the optimal policy in the third and fourth fraction is analogous to the second fraction. The dose delivered in the last fraction corresponds to the residual BED that can be delivered to the OAR to meet the constraint of 90 Gy BED3. In figure 3(b) of the original paper, the value function of the first fraction is illustrated. For a sparing factor $\delta_1 = 0.99$, a tumor BED10 of 53 Gy is expected based on the initial probability distribution.

Figure 3 of the original paper shows the optimal policy that was calculated based on the initial estimate of the probability distribution, $P(\delta; \mu_1, \sigma_1)$. By updating the probability distribution in each fraction as described in section~2.7, the additional sparing factor observations can be incorporated into a better estimate of the patient's geometric variability. In each fraction, the optimal policy is recalculated based on the updated probability distribution and followed for the current fraction. For example, in fraction two the sparing factor $\delta_2 = 0.87$ is observed. As a result, the probability distribution is updated to $P(\delta; \mu_2, \sigma_2)$ with $\mu_2 = 0.91$ and $\sigma_2 = 0.057$. The reoptimized policy proposes a dose of 8.8 Gy in fraction two. For patient 7, the change in the probability distribution between fractions is small and consequently has only a small impact on the optimal policy, which is illustrated in figure 9 in appendix A.3 of the original paper.

Table~\ref{tab:patient7} summarizes the doses that would be delivered in each fraction and compares AF to the reference treatment and the upper bound. In the reference treatment, the doses delivered to the tumor are $6/\delta_t$, resulting in a cumulative tumor BED10 of 50.0 Gy. AF increases the tumor BED10 to 50.2 Gy, mainly by delivering a larger dose in fraction 2, exploiting the lower sparing factor. At the beginning of the treatment, the achievable tumor BED was estimated to be 53 Gy BED10. The resulting treatment plan delivered less BED to the tumor, as the observed sparing factors were mostly higher than what was expected based on the first two sparing factors.

A comparison to the upper bound indicates that, for the sparing factor variations observed in this patient, the possible improvement that AF may achieve is limited. By knowing all sparing factors in advance, one could have further increased the dose delivered in fraction 2. However, even in this hypothetical case the tumor BED10 increases by only 1.7 Gy.

\subsection{Evaluation of AF for all patients}

The AF algorithm was applied to all 16 patients. Sparing factors and corresponding doses delivered in each fraction are shown in figure 4 of the original paper and reported in table 3 in appendix A.4. The quantitative comparison to the reference treatment is provided in table~\ref{tab:all-patients}.

\begin{table}[ht]
\centering
\caption{Comparison of treatment plans in Gy BED10. The difference column is calculated as AF BED minus reference plan BED, i.e. positive values indicate an improvement using AF.}
\label{tab:all-patients}
\begin{tabular}{rrrrr}
\toprule
Patient & Reference plan BED & Upper bound & Adaptive fractionation & Difference \\
\midrule
1 & 49.1 & 49.4 & 49.2 & 0.1 \\
2 & 55.6 & 55.7 & 55.6 & 0.1 \\
3 & 81.0 & 108.7 & 100.6 & 19.6 \\
4 & 49.7 & 50.6 & 50.2 & 0.5 \\
5 & 49.7 & 53.9 & 51.3 & 1.6 \\
6 & 43.4 & 43.8 & 43.5 & 0.1 \\
7 & 50.0 & 51.7 & 50.2 & 0.2 \\
8 & 71.3 & 93.1 & 85.5 & 14.2 \\
9 & 46.0 & 46.3 & 46.2 & 0.2 \\
10 & 67.3 & 68.8 & 67.5 & 0.2 \\
11 & 61.2 & 62.0 & 61.4 & 0.2 \\
12 & 54.1 & 55.3 & 54.9 & 0.8 \\
13 & 69.3 & 108.4 & 64.1 & -5.2 \\
14 & 63.0 & 63.6 & 63.3 & 0.3 \\
15 & 69.1 & 70.9 & 69.7 & 0.6 \\
16 & 52.5 & 63.0 & 51.5 & -1.0 \\
\bottomrule
\end{tabular}
\end{table}

For 14 out of the 16 patients, AF yields an increase in tumor BED10 compared to the reference treatment. However, for most patients, the difference is only around one Gray or less. For patients with small variations in the sparing factor, such as patients 1, 2, and 9, there is little variation in the dose per fraction. Consequently, AF performs similar to the reference treatment. In fact, the calculation of the upper bound shows that the benefit of AF is a priori limited to less than 0.3 Gy for these three patients. Patient 5 shows an intermediate variability in the sparing factor. Here, AF realizes 1.6 Gy tumor BED10 increase, which corresponds to approximately half of the upper bound for the improvement.

Patients 3, 8, and 13 show large variations in the sparing factor and larger differences between AF and the reference treatment. For patient 8, a sparing factor of 0.58 is observed in fraction 3, which is substantially lower compared to what was observed before. This is exploited by delivering a large dose of 20.2 Gy, which corresponds to most of the residual BED3 that is allowed in the OAR. In fractions 4 and 5, the sparing factor is again higher. Thus, delivering a large dose in fraction 3 was indeed a good decision, resulting in an improvement of 14.2 Gy tumor BED10 compared to the reference treatment.\footnote{The impact of enforcing constraints on the minimum and maximum dose delivered to the tumor in each fraction is discussed in section~4.1.} Similarly, patient 3 had a very low sparing factor in fraction 3, leading to a large tumor BED10 increase of 19.6 Gy.

The clinically delivered treatment for patient 8 was based on a fixed prescription of 7 Gy per fraction to the PTV and an OAR constraint of 6 Gy per fraction. The DVHs are shown in figure 8(b) in appendix A.2 of the original paper. In the clinical treatment, the lower sparing factor in fraction 3 translates into a lower OAR dose rather than an increased tumor dose. As a consequence, the accumulated BED3 in the OAR is only 76.4 Gy and stays below the limit of 90 Gy. A cumulative BED10 of 59.2 Gy was delivered to the PTV.

\begin{figure}[ht]
\centering
\fbox{\parbox{0.8\linewidth}{Figure omitted in this text transcription. See the original PDF for the rendered graphic.}}
\caption{Optimal policy for fractions one to five assuming $P(\delta; \mu_1, \sigma_1)$ with $\mu_1 = 0.94$ and $\sigma_1 = 0.058$. Panel (b) shows the value function of the first fraction. The crosses mark the state of the treatment in each fraction, given the observed daily sparing factors and delivered doses for patient 7.}
\label{fig:policy-patient7}
\end{figure}

\begin{figure}[ht]
\centering
\fbox{\parbox{0.8\linewidth}{Figure omitted in this text transcription. See the original PDF for the rendered graphic.}}
\caption{Sparing factors and corresponding tumor doses for each fraction for all patients. Different fractions are color coded; sparing factors are shown as crosses according to the axis on the right; delivered doses are marked with stars according to the axis on the left.}
\label{fig:all-patient-doses}
\end{figure}

\subsection{Dependence on the order of sparing factors}

Whereas patients 3 and 8 benefit from the large variation in sparing factors through AF, we observe for patient 13 that the reference treatment performs better. This can be explained through the order of sparing factors, $[1.06, 0.92, 0.84, 0.82, 1.01, 0.53]$, which is unfavorable for patient 13. The highest sparing factor is observed for the planning MR and the lowest sparing factor is observed in the last fraction. The dose delivered in the last fraction simply corresponds to the residual BED3 that can be delivered to the OAR. Since the algorithm cannot anticipate an exceptionally low sparing factor in the last fraction, it does not lower the dose in previous fractions to exploit the low sparing factor in the last fraction. In addition, observing the highest sparing factor in the planning MR leads to relatively large doses delivered in fractions 1 and 2. As a consequence, a residual dose of only 4.2 Gy can be delivered to the OAR in the last fraction, resulting in worse performance than the reference treatment.

The impact of the sparing factor order on the performance of AF is further analysed in figure 5 of the original paper for patients 8 and 13. We consider permutations of the sparing factors, investigating hypothetical treatments in which the same values were observed but in different order. When the lowest sparing factor is observed in the last fraction, AF will in most cases yield lower cumulative tumor BED10 compared to the reference treatment. Instead, if the lowest sparing factor is observed in the second or third fraction, AF improves on the reference treatment.\footnote{Note that permuting the sparing factors only affects the performance of AF whereas the cumulative tumor BED10 of the reference treatment remains the same.}

\begin{figure}[ht]
\centering
\fbox{\parbox{0.8\linewidth}{Figure omitted in this text transcription. See the original PDF for the rendered graphic.}}
\caption{Histogram of differences in BED10 (adaptive fractionation minus reference treatment), for all 720 permutations of the sparing factors for (a) patient 8, $[0.77, 0.88, 0.8, 0.58, 0.86]$ and (b) patient 13, $[1.06, 0.92, 0.84, 0.82, 1.01, 0.53]$. Highlighted are the 120 permutations for which the lowest sparing factor was observed in the last fraction (red), the planning session (blue), and the second fraction (green).}
\label{fig:permutations}
\end{figure}

\subsection{Simulated patients}

To quantify the benefit of AF for a larger patient cohort, additional patients were generated by randomly drawing sparing factors from a predefined distribution. First, a patient-specific mean $\mu_p$ and standard deviation $\sigma_p$ was selected for each generated patient $p$. The standard deviations $\sigma_p$ were drawn from the gamma distribution illustrated in figure 2 such that the variation of the sampled patients resembles the observed variation in our cohort of 16 patients. Similarly, the patient-specific means $\mu_p$ were drawn from a normal distribution that was estimated based on the means of the 16 extracted patients. Subsequently, 6 sparing factors were drawn from a normal distribution with the corresponding $\mu_p$ and $\sigma_p$.

Figure 6 of the original paper shows the histogram of tumor BED10 differences between AF and the reference treatment for 5000 generated patients. The mean benefit of AF for this patient cohort is 0.93 Gy BED10. 81.2\% of the sampled patients had a better treatment when AF was applied. The histogram shows a long tail in the positive direction, corresponding to the relatively few patients that have a large benefit from AF. The less extended tail in the negative direction shows that it is unlikely but possible that AF leads to lower tumor BED10.

To study the dependence of the benefit of AF on the amount of variation of the sparing factors, patients from 11 different populations have been sampled. All populations have a constant mean sparing factor $\mu = 0.9$; the standard deviation is constant within each population but differs between populations. For each population, the optimal policy was calculated assuming that the probability distribution is known a priori. Thus, the results are slightly superior compared to the situation that the probability has to be estimated from the observed sparing factors. Figure 7 of the original paper shows the distribution of BED10 difference between AF and the reference treatment as a function of the sparing factor standard deviation. Larger variation in the sparing factors is clearly correlated with a larger mean benefit from AF. Figure 7 also shows that the spread of BED10 differences is larger for increasing standard deviations, i.e. more extreme treatments occur. Large benefits from AF are more likely than treatments that are substantially worse than the reference treatment, explaining the increase in the mean benefit.

\begin{figure}[ht]
\centering
\fbox{\parbox{0.8\linewidth}{Figure omitted in this text transcription. See the original PDF for the rendered graphic.}}
\caption{Histogram of the difference in tumor BED10 in Gy between adaptive fractionation and reference treatment for 5000 sampled patients.}
\label{fig:sampled-patients}
\end{figure}

\begin{figure}[ht]
\centering
\fbox{\parbox{0.8\linewidth}{Figure omitted in this text transcription. See the original PDF for the rendered graphic.}}
\caption{Box plot of BED10 differences between AF and the reference treatment. The larger the standard deviation of the sparing factor distribution, the larger the spread of the BED10 differences. The red line shows the median, blue lines the 25\% and 75\% percentiles, black lines a 1.5-fold extension of the interquartile range.}
\label{fig:std-dependence}
\end{figure}

\section{Modifications and Extensions}

The approach to AF described above can be modified and extended in different ways. In this section, selected modifications are described.

\subsection{Constraints on the dose per fraction}

In the results above we did not consider a minimum tumor dose that must be delivered in each fraction. In addition, the maximum dose was allowed to be high. In practice, one may want to limit the dose per fraction $d_t$ to be within a range of clinically applied fractionation schemes. As an example, we consider a minimum dose of 3 Gy and a maximum dose of 15 Gy per fraction. For most of the 16 patients, this would lead to only minor changes in the treatment since the optimal doses are within this range anyway (figure 4 in the original paper). For patients with large sparing factor variations, modifications are seen. For patient 8, doses of $[4.1, 8.2, 15, 3.4, 8.7]$ would be delivered, resulting in a cumulative tumor BED10 of 78.9 Gy. Thus, constraints on the dose per fraction reduce the potential benefit of AF for such patients. For patient 13 these constraints on the dose per fraction result in an increase of 0.5 Gy BED10 compared to the unconstrained AF plan due to a larger dose delivered in the last fraction. Hence, the minimum dose also reduces the potential for outcomes which are substantially inferior compared to the reference plan.

\subsection{Alternative definitions of the sparing factor}

In this work, the sparing factor was defined as the dose exceeded in 1 cc of the relevant OAR divided by the dose exceeded in 95\% of the PTV, which is motivated by two considerations: first, clinical planning and reporting is based on these parameters, and second, it can be assumed that the sparing factor could not have been improved in the treatment plan optimization step. However, generally the definition of the sparing factor is ambiguous. For example, the sparing factor can be defined using the dose exceeded in 95\% of the GTV rather than the PTV. This would change the results of adaptive fractionation quantitatively but not qualitatively. A comparison is provided in appendix A.5 in table 4.

\subsection{Minimizing OAR BED for a given tumor prescription}

In this paper, we considered patients in whom tumor BED was compromised due to the proximity to an OAR, and thus the goal was to maximize cumulative tumor BED while respecting an OAR constraint. For patients with lower sparing factors, a desired tumor prescription $B_T^{\min}$ may be achieved. In this case, the objective may become the minimization of OAR BED rather than a further increase of tumor BED. To account for that, the MDP model can be extended to a three-dimensional state

\begin{equation}
s_t = \left(
\delta_t,
\sum_{\tau=1}^{t-1}\left(\delta_{\tau} d_{\tau} + \frac{\delta_{\tau}^2 d_{\tau}^2}{(\alpha/\beta)_N}\right),
\sum_{\tau=1}^{t-1}\left(d_{\tau} + \frac{d_{\tau}^2}{(\alpha/\beta)_T}\right)
\right),
\label{eq:state-three-dim}
\end{equation}

where the third component is the tumor BED10 accumulated up to fraction $t$. The immediate cost $r_t$ in fraction $t$ is given by the BED delivered to the OAR in that fraction

\begin{equation}
r_t = -\left(\delta_t d_t + \frac{\delta_t^2 d_t^2}{(\alpha/\beta)_N}\right),
\label{eq:oar-cost}
\end{equation}

and the terminal cost is

\[
v_{F+1}(\delta, B_N, B_T) =
\begin{cases}
-\infty & \text{if } B_N > B_N^{\max}, \\
-\kappa \left(B_T^{\min} - B_T\right) & \text{if } B_N \leq B_N^{\max} \text{ and } B_T < B_T^{\min}, \\
0 & \text{otherwise.}
\end{cases}
\]

The parameter $\kappa$ weights the cost for underdosing the tumor. As long as $\kappa$ is chosen large enough, the optimal policy will prioritize delivering $B_T^{\min}$ to the tumor over the objective of minimizing OAR BED.

As an example, we assume a prescription BED10 of $B_T^{\min} = 72$ Gy, corresponding to 5 fractions of 8 Gy. For the patients in whom the achievable BED10 is substantially lower, this extended DP algorithm yields almost identical results as reported in figure 4 of the original paper. This is because, in the range of relevant values of $\delta_t$ and $B_N$, the optimal policy is almost independent of $B_T$ and approximately equal to the algorithm described in section~2.5. However, for patient 8, the algorithm would deliver doses of $[4.1, 8.3, 17.4, 0.9, 2.0]$ Gy, leading to a tumor BED10 of 72 Gy and an OAR BED3 of 76.6 Gy.

\section{Discussion}

\subsection{Summary of main findings}

For SBRT treatments of abdominal lesions at the MR-linac, target coverage may be compromised when tumors are located close to an OAR. In this work, we investigate whether AF may improve tumor BED. In current practice, treatments are typically based on a fixed prescription dose for each fraction. For tumors that are close to an OAR in the planning MR scan, the prescription dose may be lowered compared to what would otherwise be desired. If in some of the fractions a larger distance of tumor and OAR would allow for a larger tumor dose, the prescription dose is not altered. Thus, the favorable geometry results in an OAR dose below the tolerance. One step toward AF, called the reference treatment in this work, consists in increasing tumor and OAR dose up to the per-fraction dose constraint of the OAR. This may come with challenges for the clinical workflow and treatment documentation but does not represent any scientific problem.

In this work, we investigate if AF can improve the ratio of tumor versus OAR BED beyond the reference treatment. To that end, the AF problem was formulated as an MDP and the optimal policy was determined via dynamic programming. In contrast to prior work, we consider the problem that the probability distribution over sparing factors is not known but has to be estimated for the individual patient. In addition, we evaluate the algorithm using real data. To that end, we analyzed 16 MR-linac patients previously treated for abdominal lesions with 5-fraction SBRT. Main findings are:

\begin{enumerate}
\item The average benefit of AF may be small. For the majority of patients that we analyzed, the tumor BED10 increase through AF was below 1 Gy. For these patients, the inter-fractional variation in geometry was too small, and calculation of an upper bound showed that no AF strategy may provide a significant improvement.
\item Significant improvements may be achieved for a subset of patients. For two patients with large variation in geometry, a substantial BED10 improvement on the order of 15 Gy could be achieved by delivering large doses when a favorable geometry occurs in the middle of the treatment.
\item Although improvements are more likely, there is a residual risk that AF yields inferior treatments. For a third patient with large geometry variation, AF resulted in 5 Gy lower tumor BED10. This problem occurs for an unfavorable order of sparing factors, e.g. if sparing factors in later fractions are substantially different from what was expected based on the initially estimated probability distribution. Presumably, no approach to AF may fully prevent this without reducing the benefit of AF for other patients.
\end{enumerate}

Simulations with randomly generated patients confirmed the main findings from the 16 analyzed patients.

\subsection{Comparison to previous publications}

Compared to the prior works by Lu et al. (2008) and Chen et al. (2008), we observe smaller benefits of AF. Both papers reported up to 30\% decrease in OAR BED using AF. The average relative difference in tumor BED in this work is on the order of a few percent. This large difference may originate from three key differences. The standard deviations of the sparing factors observed for our MR-linac patients were notably lower than the standard deviations assumed in previous papers, where standard deviations between 0.1 and 0.6 were used to model the variability of the sparing factor. Furthermore, the number of fractions is larger in both earlier papers, where the prescribed dose to the tumor was delivered in 40 fractions, which gives more opportunity to the algorithm to deliver lower doses on bad days and higher doses on good days. Additionally, in both earlier papers the probability distribution of the sparing factors was assumed to be known a priori, which overestimates the quality of the treatments compared to the real-world situation where the distribution has to be estimated for the individual patient during treatment.

\subsection{Coupling of treatment planning and fractionation decision}

In this work, the AF problem was decoupled from the treatment planning problem. We used the clinically delivered plans that were created based on fixed prescription doses and OAR constraints, and assumed that AF is performed by upscaling or downscaling these treatment plans. In principle, one may expect improvements by considering both problems jointly. This could well be done for the last fraction by setting the residual dose that may be delivered to the OAR as the constraint. However, it is unclear how to combine treatment planning and fractionation decision for earlier fractions.

\section{Conclusion}

Image guidance and daily replanning at the MR-linac enables a clinical implementation of AF as an approach to exploit day-to-day variations in the distance of the tumor from the dose-limiting OAR. Based on our study considering 5-fraction SBRT treatments of abdominal lesions in proximity to bowel, stomach, or duodenum, we conclude that for the majority of patients, the amount of interfraction motion may be too small to substantially benefit from AF. However, substantial tumor dose escalation may be achieved for a subset of patients with large day-to-day changes of the geometry.

\appendix
\section{Appendix}

\subsection{GitHub repository}

The algorithm and a graphical interface to apply AF are available in the GitHub repository \texttt{adaptfx}. The version used for the results in this publication is made available as a persistent release:
\url{https://github.com/openAFT/adaptfx/tree/perez_haas}. This release also contains the DVHs of all 16 patients inside CSV files.

\subsection{DVH of clinically delivered treatments}

\begin{figure}[ht]
\centering
\fbox{\parbox{0.8\linewidth}{Figure omitted in this text transcription. See the original PDF for the rendered graphic.}}
\caption{DVH of GTV, PTV, and the respective dose-limiting OAR for selected patients for the initial plan and the five fractions. For the OAR, only the 5 cc receiving the highest dose is included in the DVH. Thus 1 cc corresponds to 20\%.}
\label{fig:dvh}
\end{figure}

\subsection{Probability update}

\begin{figure}[ht]
\centering
\fbox{\parbox{0.8\linewidth}{Figure omitted in this text transcription. See the original PDF for the rendered graphic.}}
\caption{(a) Update of the probability distribution of the sparing factors for patient 7 between fraction 1 and fraction 4. (b) Impact of the probability distribution update on the optimal policy. Shown is the optimal dose to deliver in fraction 4 as a function of the sparing factor for a previously delivered OAR BED3 given by 47.6, showing differences of up to approximately 1 Gy.}
\label{fig:prob-update}
\end{figure}

\subsection{Sparing factors and doses per fraction}

\setlength{\LTleft}{0pt}
\setlength{\LTright}{0pt}
\small
\begin{longtable}{llllllll}
\caption{Delivered doses in adaptive fractionation and the treatment corresponding to the upper bound.}
\label{tab:appendix-doses}\\
\toprule
Patient & Plan type & Planning & Fraction 1 & Fraction 2 & Fraction 3 & Fraction 4 & Fraction 5 \\
\midrule
\endfirsthead
\toprule
Patient & Plan type & Planning & Fraction 1 & Fraction 2 & Fraction 3 & Fraction 4 & Fraction 5 \\
\midrule
\endhead
1 & sparing factor & 0.99 & 0.95 & 0.98 & 0.96 & 1.02 & 1.01 \\
  & adaptive & --- & 6.8 & 5.7 & 6.5 & 5.0 & 6.5 \\
  & upper bound & --- & 7.7 & 6.0 & 7.1 & 4.6 & 4.9 \\
\midrule
2 & sparing factor & 0.90 & 0.91 & 0.92 & 0.88 & 0.90 & 0.89 \\
  & adaptive & --- & 6.3 & 6.2 & 7.8 & 6.5 & 6.5 \\
  & upper bound & --- & 6.0 & 5.6 & 7.9 & 6.6 & 7.2 \\
\midrule
3 & sparing factor & 0.73 & 0.67 & 0.78 & 0.53 & 0.89 & 0.73 \\
  & adaptive & --- & 10.8 & 3.9 & 21.8 & 0.3 & 2.4 \\
  & upper bound & --- & 1.6 & 0.4 & 27.9 & 0.0 & 0.8 \\
\midrule
4 & sparing factor & 0.95 & 1.00 & 0.92 & 0.92 & 1.04 & 1.01 \\
  & adaptive & --- & 5.1 & 7.9 & 7.0 & 4.0 & 6.5 \\
  & upper bound & --- & 4.7 & 8.7 & 8.7 & 3.7 & 4.4 \\
\midrule
5 & sparing factor & 0.94 & 0.83 & 0.97 & 1.13 & 1.05 & 0.96 \\
  & adaptive & --- & 9.3 & 3.9 & 2.4 & 5.1 & 9.0 \\
  & upper bound & --- & 14.9 & 4.0 & 1.7 & 2.5 & 4.2 \\
\midrule
6 & sparing factor & 1.04 & 1.03 & 1.10 & 1.04 & 1.14 & 1.08 \\
  & adaptive & --- & 5.9 & 4.5 & 6.3 & 4.4 & 6.6 \\
  & upper bound & --- & 7.3 & 4.6 & 6.7 & 3.8 & 5.2 \\
\midrule
7 & sparing factor & 0.88 & 0.99 & 0.87 & 0.98 & 1.04 & 1.00 \\
  & adaptive & --- & 4.4 & 8.8 & 4.5 & 4.4 & 8.2 \\
  & upper bound & --- & 4.4 & 12.6 & 4.7 & 3.3 & 4.1 \\
\midrule
8 & sparing factor & 0.77 & 0.88 & 0.80 & 0.58 & 0.86 & 0.77 \\
  & adaptive & --- & 4.0 & 8.2 & 20.2 & 1.0 & 2.3 \\
  & upper bound & --- & 0.49 & 1.1 & 25.3 & 0.6 & 1.4 \\
\midrule
9 & sparing factor & 1.04 & 1.02 & 0.99 & 1.02 & 1.08 & 1.05 \\
  & adaptive & --- & 6.1 & 6.7 & 5.5 & 4.6 & 6.1 \\
  & upper bound & --- & 6.0 & 7.6 & 6.0 & 4.2 & 5.0 \\
\midrule
10 & sparing factor & 0.86 & 0.78 & 0.79 & 0.81 & 0.83 & 0.73 \\
   & adaptive & --- & 9.4 & 7.7 & 6.5 & 6.2 & 8.2 \\
   & upper bound & --- & 6.8 & 6.1 & 5.0 & 4.4 & 14.0 \\
\midrule
11 & sparing factor & 1.01 & 0.83 & 0.80 & 0.82 & 0.85 & 0.91 \\
   & adaptive & --- & 10.6 & 8.7 & 6.0 & 4.9 & 4.4 \\
   & upper bound & --- & 7.2 & 10.2 & 8.0 & 5.9 & 4.0 \\
\midrule
12 & sparing factor & 0.88 & 0.94 & 0.88 & 0.97 & 0.96 & 0.85 \\
   & adaptive & --- & 5.1 & 7.7 & 4.7 & 6.1 & 8.9 \\
   & upper bound & --- & 4.8 & 7.9 & 3.9 & 4.2 & 10.9 \\
\midrule
13 & sparing factor & 1.06 & 0.92 & 0.84 & 0.82 & 1.01 & 0.53 \\
   & adaptive & --- & 8.5 & 9.6 & 7.1 & 2.6 & 7.9 \\
   & upper bound & --- & 0.0 & 0.1 & 0.2 & 0.0 & 28.3 \\
\midrule
14 & sparing factor & 0.84 & 0.85 & 0.83 & 0.78 & 0.82 & 0.84 \\
   & adaptive & --- & 6.7 & 7.8 & 9.7 & 6.1 & 5.9 \\
   & upper bound & --- & 5.2 & 6.3 & 11.4 & 7.0 & 5.7 \\
\midrule
15 & sparing factor & 0.78 & 0.72 & 0.74 & 0.83 & 0.79 & 0.79 \\
   & adaptive & --- & 9.9 & 7.7 & 4.1 & 7.3 & 9.3 \\
   & upper bound & --- & 13.3 & 9.6 & 3.7 & 5.3 & 5.3 \\
\midrule
16 & sparing factor & 0.98 & 1.05 & 1.06 & 1.06 & 0.90 & 0.74 \\
   & adaptive & --- & 4.7 & 5.0 & 5.5 & 9.8 & 5.8 \\
   & upper bound & --- & 1.3 & 1.3 & 1.3 & 3.0 & 18.8 \\
\bottomrule
\end{longtable}
\normalsize

\begin{figure}[ht]
\centering
\fbox{\parbox{0.8\linewidth}{Figure omitted in this text transcription. See the original PDF for the rendered graphic.}}
\caption{PTV and the corresponding GTV sparing factors of all patients.}
\label{fig:ptv-gtv}
\end{figure}

\subsection{Results for GTV-based sparing factors}

The same analysis as described in sections 2 and 3 can be conducted with GTV-based sparing factors, i.e. all sparing factors are defined via the $D_{95}$ of the GTV rather than the PTV, and the goal is to maximize BED10 in the GTV. Figure~\ref{fig:ptv-gtv} shows the relation of GTV and PTV sparing factors. As expected, GTV and PTV sparing factors are correlated. Consequently, AF yields similar treatments for both definitions of the sparing factor. Table~\ref{tab:gtv-ptv} (column 1) reports BED10 differences in the GTV compared to the reference treatment for AF optimized for GTV-based sparing factors. The same treatment can then be evaluated for BED10 in the PTV (column 2). For comparison, table~\ref{tab:gtv-ptv} also reports BED10 differences in the PTV (column 3) and the GTV (column 4) for AF optimized for PTV-based sparing factors.

\begin{table}[ht]
\centering
\caption{Plan differences based on GTV optimization (columns one and two) and PTV optimization (columns three and four). The optimal doses in the first two columns have been computed based on the GTV sparing factors and based on the PTV sparing factors in columns three and four. The respective differences to the reference plans are given in BED10.}
\label{tab:gtv-ptv}
\begin{tabular}{rrrrr}
\toprule
Patient & GTV diff. (GTV opt.) & PTV diff. (GTV opt.) & PTV diff. (PTV opt.) & GTV diff. (PTV opt.) \\
\midrule
1 & 0.4 & -0.2 & 0.1 & 0.3 \\
2 & 0.0 & 0.0 & 0.0 & 0.0 \\
3 & 21.8 & 11.7 & 19.6 & 36.0 \\
4 & 4.6 & 0.2 & 0.5 & 2.3 \\
5 & 1.5 & 1.8 & 1.6 & 0.6 \\
6 & 0.6 & -0.2 & 0.1 & -0.8 \\
7 & -0.8 & -0.4 & 0.2 & -0.2 \\
8 & 17.7 & 13.6 & 14.2 & 18.6 \\
9 & -2.6 & -3.1 & 0.2 & 0.7 \\
10 & 1.0 & 0.3 & 0.2 & 0.6 \\
11 & -0.4 & 0.0 & 0.1 & -0.2 \\
12 & 1.5 & 1.0 & 0.8 & 0.9 \\
13 & 13.8 & 8.2 & -5.2 & -6.9 \\
14 & -2.1 & -0.7 & 0.3 & 0.4 \\
15 & 0.7 & -0.2 & 0.6 & 1.3 \\
16 & 2.8 & -0.3 & -1.0 & -0.2 \\
Mean & 3.8 & 2.0 & 2.0 & 3.4 \\
\bottomrule
\end{tabular}
\end{table}

Comparing results based on GTV and PTV-based sparing factors, the most notable difference is seen for patient 13, where AF based on PTV sparing factors leads to lower BED10 than the reference treatment. For GTV-based sparing factors, we instead see an improvement through AF, even though the lowest sparing factor is still observed in the last fraction. However, there is overall less variation in the GTV sparing factors up to fraction 4, the initial sparing factor better represents the mean, and the sparing factor in fraction 4 was higher than the other sparing factors, leading to a small dose to be delivered in fraction 4. In combination, this results in a larger residual dose available in fraction 5.

\section*{ORCID iDs}

Y. P\'erez Haas: \url{https://orcid.org/0000-0002-4282-0643}\\
R. Ludwig: \url{https://orcid.org/0000-0001-9434-328X}\\
R. Dal Bello: \url{https://orcid.org/0000-0002-8755-377X}

\begin{thebibliography}{99}

\bibitem{ref1} Acharya S et al. 2016. Online magnetic resonance image guided adaptive radiation therapy: first clinical applications. \textit{International Journal of Radiation Oncology, Biology, Physics} 94, 394--403.
\bibitem{ref2} Andratschke N et al. 2018. The SBRT database initiative of the German Society for Radiation Oncology (DEGRO): patterns of care and outcome analysis of stereotactic body radiotherapy (SBRT) for liver oligometastases in 474 patients with 623 metastases. \textit{BMC Cancer} 18, 283.
\bibitem{ref3} Brock K K. 2019. Adaptive radiotherapy: moving into the future. \textit{Seminars in Radiation Oncology} 29, 181--184.
\bibitem{ref4} Chen M, Lu W, Chen Q, Ruchala K and Olivera G. 2008. Adaptive fractionation therapy: II. Biological effective dose. \textit{Physics in Medicine \& Biology} 53, 5513--5525.
\bibitem{ref5} Chen W, Gemmel A and Rietzel E. 2013. A patient-specific planning target volume used in `plan of the day' adaptation for interfractional motion mitigation. \textit{Journal of Radiation Research} 54, i82--i90.
\bibitem{ref6} Fowler J F. 2006. Development of radiobiology for oncology: a personal view. \textit{Physics in Medicine \& Biology} 51.
\bibitem{ref7} Guckenberger M, Wilbert J, Richter A, Baier K and Flentje M. 2011. Potential of adaptive radiotherapy to escalate the radiation dose in combined radiochemotherapy for locally advanced non-small cell lung cancer. \textit{International Journal of Radiation Oncology, Biology, Physics} 79.
\bibitem{ref8} Henke L et al. 2018. Phase I trial of stereotactic MR-guided online adaptive radiation therapy (SMART) for the treatment of oligometastatic or unresectable primary malignancies of the abdomen. \textit{Radiotherapy and Oncology} 126, 519--526.
\bibitem{ref9} ICRU. 1993. Prescribing, recording, and reporting photon beam therapy. \textit{ICRU Report 50}, 3--26.
\bibitem{ref10} ICRU. 1999. Prescribing, recording, and reporting photon beam therapy (supplement to ICRU Report 50). \textit{ICRU Report 62}, 3--20.
\bibitem{ref11} Jones B, Dale R G, Deehan C, Hopkins K I and Morgan D A L. 2001. The role of biologically effective dose (BED) in clinical oncology. \textit{Clinical Oncology} 13, 71--81.
\bibitem{ref12} Kl\"uter S. 2019. Technical design and concept of a 0.35 T MR-linac. \textit{Clinical and Translational Radiation Oncology} 18, 98--101.
\bibitem{ref13} Lajtha L G, Oliver R and Ellis F. 1960. Rationalisation of fractionation in radiotherapy. \textit{British Journal of Radiology} 33, 634--635.
\bibitem{ref14} Lo S S et al. 2010. Stereotactic body radiation therapy: a novel treatment modality. \textit{Nature Reviews Clinical Oncology} 7, 44--54.
\bibitem{ref15} Lu W, Chen M, Chen Q, Ruchala K and Olivera G. 2008. Adaptive fractionation therapy: I. Basic concept and strategy. \textit{Physics in Medicine \& Biology} 53, 5495--5511.
\bibitem{ref16} Mayinger M et al. 2021. Benefit of replanning in MR-guided online adaptive radiation therapy in the treatment of liver metastasis. \textit{Radiation Oncology} 16, 84.
\bibitem{ref17} McMahon S J. 2018. The linear quadratic model: usage, interpretation and challenges. \textit{Physics in Medicine \& Biology} 64, 01TR01.
\bibitem{ref18} Palacios M A et al. 2018. Role of daily plan adaptation in MR-guided stereotactic ablative radiation therapy for adrenal metastases. \textit{International Journal of Radiation Oncology, Biology, Physics} 102, 426--433.
\bibitem{ref19} Pavic M et al. 2022. MR-guided adaptive stereotactic body radiotherapy (SBRT) of primary tumor for pain control in metastatic pancreatic ductal adenocarcinoma (mPDAC): an open randomized, multicentric, parallel group clinical trial (MASPAC). \textit{Radiation Oncology} 17.
\bibitem{ref20} Ramakrishnan J, Craft D, Bortfeld T and Tsitsiklis J N. 2012. A dynamic programming approach to adaptive fractionation. \textit{Physics in Medicine \& Biology} 57, 1203--1216.
\bibitem{ref21} Sutton R S and Barto A G. 2018. \textit{Reinforcement Learning: An Introduction}. 2nd edn. Cambridge, MA: MIT Press.
\bibitem{ref22} Tyagi N et al. 2021. Feasibility of ablative stereotactic body radiation therapy of pancreas cancer patients on a 1.5 tesla magnetic resonance-linac system using abdominal compression. \textit{Physics and Imaging in Radiation Oncology} 19, 53--59.
\bibitem{ref23} van Herk M, Remeijer P, Rasch C and Lebesque J V. 2000. The probability of correct target dosage: dose-population histograms for deriving treatment margins in radiotherapy. \textit{International Journal of Radiation Oncology, Biology, Physics} 47, 1121--1135.
\bibitem{ref24} van Leeuwen C M et al. 2018. The alpha and beta of tumours: a review of parameters of the linear-quadratic model, derived from clinical radiotherapy studies. \textit{Radiation Oncology} 13, 96.
\bibitem{ref25} Wu C, Jeraj R, Olivera G H and Mackie T R. 2002. Re-optimization in adaptive radiotherapy. \textit{Physics in Medicine \& Biology} 47, 3181--3195.
\bibitem{ref26} Wulf J et al. 2006. Stereotactic radiotherapy of primary liver cancer and hepatic metastases. \textit{Acta Oncologica} 45, 838--847.

\end{thebibliography}

\end{document}
